%=======================================================================
%  CHAPTER 4 — MODULATION AND DEMODULATION METHODS  (4 hrs)
%=======================================================================
\section{Modulation--Demodulation Methods in Mobile Communications}

{\color{sub}\footnotesize 4 hrs \gap$\cdot$\gap asked in \mk{20 of 22 papers}
\gap$\cdot$\gap 5 syllabus sub-topics \gap$\cdot$\gap
\mk{no paper has ever set a numerical here} (one worked example is given anyway, \S4.4).}

\lead{Read this first:}
\begin{itemize}
  \item \mk{This is a drawing chapter.} Almost every question says +Fig, block diagram or
        constellation diagram. Marks come from the \textbf{diagram}, not the prose.
  \item Four blocks carry the chapter, in order of return:
  \begin{itemize}
    \item \textbf{OFDM} \tS{9} $\rightarrow$ transmitter + receiver block diagram, always 4 or 8 marks.
    \item \textbf{Spread spectrum, DSSS and FHSS} \tS{7} $\rightarrow$ definition $+$ block diagram.
    \item \textbf{MSK / GMSK} \tS{6} $\rightarrow$ why GSM picked GMSK.
    \item \textbf{QPSK} \tF{4} $\rightarrow$ transmitter, receiver, constellation.
  \end{itemize}
  \item Everything else (BPSK, DPSK, OQPSK, $\pi/4$-QPSK, QAM) is a \textbf{one-paper}
        question, but each was worth 6 to 8 marks when it came.
  \item \mk{Learn 3 diagrams cold}: QPSK transmitter/receiver pair, OFDM transmitter/receiver
        pair, DSSS + FHSS transmitter pair. They answer roughly 30 marks of the paper.
  \item FHSS appears here as a \emph{modulation} and again in Ch 7 as \textbf{FHMA}, a
        \emph{multiple access} method. Same block diagram, different framing. Draw it once,
        label it for whichever chapter is asked.
\end{itemize}

\hr

\T{4.1 Why modulate, why digital, and how schemes are judged}

{\color{sub}\footnotesize Only one PYQ lives in this section, and it is worth 2 marks. Everything
else here is \textbf{vocabulary the rest of the chapter assumes}: $\eta_p$, $\eta_B$, the
trade-off between them, and what a constellation diagram actually plots. Skim it once, then
never look at it again.}

\Q{Advantages of digital modulation over analog}{\m{2}}
{\color{sub}\footnotesize \yr{\textbf{\texttt{79 Bh}} \m{2}} \gap$\cdot$\gap
\mk{always a rider on the GMSK question, never asked alone}}

\sbs{%
\lead{Modulation itself}
\begin{itemize}
  \item Process of impressing message on high-freq carrier suitable for transmission.
  \item Translates baseband (modulating) signal $\rightarrow$ bandpass (modulated) signal.
  \item \textbf{Why:} antenna size $\propto \lambda$, so baseband speech at 3 kHz needs
        km-long antennas.
  \item \textbf{Why:} shifting to different carriers lets many users share one medium
        (multiple access).
  \item \textbf{Why:} higher freq $\rightarrow$ longer range, smaller hardware.
\end{itemize}}{%
\lead{Digital over analog \m{2}}
\begin{itemize}
  \item Regenerable: noise does not accumulate, receiver rebuilds clean symbols.
  \item Resistant to fading, noise, interference.
  \item Voice $+$ data $+$ video multiplexed on one channel.
  \item Encryption possible $\rightarrow$ security and privacy.
  \item Error control coding detects and corrects errors (Ch 6).
  \item Built in DSP software, not analog hardware $\rightarrow$ cheap, small, reprogrammable.
  \item Equalization and diversity apply cleanly (Ch 5).
  \item \mk{Cost:} needs more bandwidth per message and needs synchronisation.
\end{itemize}}

\penalty0\vspace{2pt}
\sbsr{0.50}{%
\lead{Power efficiency $\eta_p$}
\begin{itemize}
  \item Ability to hold message quality at low power (low SNR).
  \item $\eta_p = E_b/N_0$ at a stated BER, usually $10^{-5}$.
  \item $E_b =$ energy per bit, $N_0 =$ noise power spectral density.
  \item Trade-off: BER $\uparrow$ as $E_b/N_0 \downarrow$.
  \item \textbf{Lower $E_b/N_0$ for the same BER $=$ more power efficient.}
\end{itemize}}{%
\lead{Bandwidth efficiency $\eta_B$}
\begin{itemize}
  \item Ability to carry data in a limited band.
  \item $\eta_B = R/B$ bps/Hz, $R =$ data rate, $B =$ RF bandwidth.
  \item Ceiling set by Shannon:
        $\eta_{B(\max)} = C/B = \log_2(1 + S/N)$.
  \item \mk{The fundamental trade-off:} improve $\eta_p$ and $\eta_B$ gets worse, and the
        reverse. No scheme wins both.
  \item M-ary buys $\eta_B$ by spending $\eta_p$.
\end{itemize}}

\penalty0\vspace{2pt}
\sbsr{0.52}{%
\lead{Choosing a scheme for a mobile link}
\begin{itemize}
  \item Low BER at low SNR (handset battery is finite).
  \item Resistance to interference (ACI, CCI) and multipath fading.
  \item Minimum occupied bandwidth (spectrum is bought, not free).
  \item Cheap, small, low-power to implement in the handset.
  \item \mk{No scheme meets all four}, so every standard is a trade-off.
\end{itemize}

\lead{Constellation (signal space) diagram}
\begin{itemize}
  \item Graphical plot of the complex envelope of every possible symbol.
  \item $x$ axis $=$ in-phase component, $y$ axis $=$ quadrature component.
  \item Distance of a point from origin $=$ amplitude, so distance$^2$ $=$ symbol power.
  \item \mk{Distance $d$ between the two nearest points sets noise immunity.}
  \item $P_e \propto Q\!\left(\dfrac{d}{\sqrt{2N_0}}\right)$, so small $d$ $\rightarrow$ high BER.
  \item Farthest point from origin $=$ peak power the amplifier must deliver.
\end{itemize}}{%
\figT{c4_dig_waveforms.png}
\penalty0\vspace{1pt}
{\color{sub}\footnotesize The three binary keying waveforms against one data stream:
(a) ASK, (b) BFSK, (c) BPSK. \yr{SST, Modulation extra note}}}

\hr

\T{4.2 Review: analog modulation and line coding \; \pillo{sub}{syllabus, no PYQ yet}}

\sbs{%
\lead{Amplitude family}
\begin{itemize}
  \item \textbf{AM (DSB-FC)}: $B = 2f_m$. Carrier carries no information, so at most
        33\,\% of power is useful. Simple envelope detector.
  \item \textbf{DSB-SC}: carrier suppressed, $B = 2f_m$. All power in sidebands, but needs
        coherent detection.
  \item \textbf{SSB}: one sideband only, $B = f_m$. \mk{Half the bandwidth, half the power},
        but sharp filters and exact carrier recovery needed.
  \item \textbf{VSB}: one full sideband $+$ vestige of the other, $B \approx 1.25 f_m$.
        Compromise used for TV video where DC content forbids true SSB.
\end{itemize}}{%
\lead{Angle family}
\begin{itemize}
  \item \textbf{FM}: instantaneous freq varies with message. Carson's rule
        $B = 2(\Delta f + f_m)$.
  \item Constant envelope $\rightarrow$ non-linear class-C amplifier usable $\rightarrow$
        power efficient.
  \item \textbf{Capture effect}: stronger of two co-channel FM signals suppresses the weaker.
  \item Wideband FM trades bandwidth for SNR $\rightarrow$ why \mk{every 1G system used FM}.
  \item \textbf{PM}: phase varies with message. FM and PM are inter-convertible through an
        integrator or differentiator.
\end{itemize}}

\penalty0\vspace{2pt}
\sbsr{0.50}{%
\lead{Line coding: what it is for}
\begin{itemize}
  \item Maps bits onto baseband pulses before modulation.
  \item Requirements: no DC content, self-clocking, spectrally compact, error detection
        possible, transparent to any bit pattern.
\end{itemize}}{%
\lead{The four to know}
\begin{itemize}
  \item \textbf{NRZ}: 1 $=$ high for full bit, 0 $=$ low. Compact, but has DC and no clock in
        long runs. \mk{GSM uses bipolar NRZ ahead of the GMSK filter.}
  \item \textbf{RZ}: pulse returns to zero mid-bit. Self-clocking, costs twice the bandwidth.
  \item \textbf{Manchester}: transition mid-bit, 1 $=$ hi$\rightarrow$lo, 0 $=$ lo$\rightarrow$hi.
        No DC, always self-clocking, twice the bandwidth. Used in 10 Mbps Ethernet.
  \item \textbf{AMI (bipolar)}: 0 $=$ zero level, 1 $=$ alternating $\pm$. No DC, and a run of
        same-polarity ones flags an error.
\end{itemize}}

\hr

\T{4.3 Linear modulation: BPSK, DPSK, QPSK, OQPSK, $\pi/4$-QPSK}

{\color{sub}\footnotesize \textbf{Linear} means the transmitted amplitude follows the
modulating signal linearly. Bandwidth efficient, but envelope varies $\rightarrow$
\mk{needs a linear (backed-off, inefficient) power amplifier}. That is exactly why mobile
handsets moved to the constant-envelope family in \S4.4.}

\Q{QPSK: transmission and detection process +Fig, constellation diagram}{\m{6}\m{7}\m{8}}
{\color{sub}\footnotesize \tF{4} \yr{\textbf{75 Bh} \m{6}, 72 Ma \m{8},
\textbf{\texttt{81 Bh}} \m{8}, \textbf{73 Bh} \m{7}} \gap$\cdot$\gap
\mk{answer $=$ two block diagrams $+$ one constellation $+$ six lines}}

\sbsr{0.47}{%
\lead{Idea}
\begin{itemize}
  \item M-ary PSK with $M = 4$, so \textbf{2 bits per symbol} (a dibit).
  \item Carrier phase takes one of 4 equally spaced values
        $\pi/4,\, 3\pi/4,\, 5\pi/4,\, 7\pi/4$.
  \item $s_i(t) = \sqrt{\dfrac{2E_s}{T_s}}\cos\!\big(2\pi f_c t + (2i-1)\tfrac{\pi}{4}\big)$,
        \; $i = 1,2,3,4$.
  \item Symbol period $T_s = 2T_b$ $\rightarrow$ \mk{half the bandwidth of BPSK for the same
        bit rate}.
  \item Null-to-null RF bandwidth $= R_b$ (BPSK needs $2R_b$), so $\eta_B = 1$ bps/Hz.
  \item \mk{BER identical to BPSK}: the dibit splits into two independent BPSK streams on
        orthogonal carriers, $P_e = Q\!\left(\sqrt{2E_b/N_0}\right)$.
  \item So QPSK gets double the bandwidth efficiency \textbf{for free}. That is the whole
        selling point.
  \item Symbols are \textbf{Gray coded} (00, 01, 11, 10 around the circle) so a wrong symbol
        costs only 1 bit.
\end{itemize}}{%
\figT{c4_qpsk_const.png}
\penalty0\vspace{1pt}
{\color{sub}\footnotesize QPSK constellation: 4 points, one per dibit, spaced $90^\circ$.
Distance to origin $= \sqrt{E_s}$, distance between neighbours $= \sqrt{2E_s} = 2\sqrt{E_b}$.
\yr{SST, Modulation extra note}}}

\penalty0\vspace{2pt}
\sbsr{0.50}{%
\figT{c4_qpsk_tx.png}
\penalty0\vspace{1pt}
{\color{sub}\footnotesize \textbf{QPSK transmitter.} \yr{SST, Fig 3.2.3}}}{%
\lead{Transmission process \m{4}}
\begin{enumerate}
  \item Binary stream at $f_b$ bps enters a \textbf{bipolar NRZ encoder} $\rightarrow$ $\pm 1$
        levels.
  \item \textbf{Serial-to-parallel converter} splits it: even-numbered bits $\rightarrow$
        $I(t)$, odd-numbered bits $\rightarrow$ $Q(t)$, each at $f_b/2$ bps.
  \item $I(t)$ multiplied by $\cos 2\pi f_c t$ in balance modulator 1.
  \item $Q(t)$ multiplied by $\sin 2\pi f_c t$ (carrier through a $90^\circ$ shifter) in
        balance modulator 2.
  \item \textbf{Summer} adds the two $\rightarrow$ QPSK signal, one of 4 phases per dibit.
  \item Output BPF limits the spectrum to the allocated channel.
\end{enumerate}}

\penalty0\vspace{2pt}
\sbsr{0.50}{%
\figT{c4_qpsk_rx.png}
\penalty0\vspace{1pt}
{\color{sub}\footnotesize \textbf{QPSK receiver}, a pair of coherent correlators.
\yr{SST, Modulation extra note}}}{%
\lead{Detection process \m{4}}
\begin{enumerate}
  \item Received QPSK splits into two arms (coherent detection, so carrier recovery circuit
        supplies $\cos$ and $\sin$ references).
  \item Upper arm: product modulator with $\cos 2\pi f_c t$, then
        \textbf{integrate-and-dump} over $0$ to $T_b$ $\rightarrow$ output $x_1$.
  \item Lower arm: product modulator with $\sin 2\pi f_c t$, then integrator $\rightarrow$
        output $x_2$.
  \item \textbf{Decision device}, threshold $0$ V:
        $x_1 > 0 \rightarrow$ 1, $x_1 < 0 \rightarrow$ 0 on the I channel, same rule for
        $x_2$ on the Q channel.
  \item \textbf{Multiplexer} interleaves the two recovered streams $\rightarrow$ original
        binary sequence.
\end{enumerate}
{\color{sub}\footnotesize \mk{Say ``coherent'' out loud.} QPSK cannot be envelope detected,
the receiver needs the carrier phase.}}

\penalty0\vspace{2pt}
\lead{The weakness that produced OQPSK and $\pi/4$-QPSK}
\begin{itemize}
  \item When \emph{both} bits of a dibit change, the phase jumps $180^\circ$ and the
        \mk{envelope passes through zero}.
  \item After the transmitter's band-limiting filter this zero crossing becomes a deep
        amplitude dip.
  \item A non-linear (efficient) amplifier hard-limits that dip $\rightarrow$ the filtered
        sidelobes \textbf{grow back} $\rightarrow$ adjacent channel interference returns.
  \item Fixing this without giving up bandwidth efficiency is the whole story of \S4.3 and
        \S4.4.
\end{itemize}

\creamq{BPSK against QPSK}{\m{5}}
{\color{sub}\footnotesize \yr{73 Ma \m{5}}}

\penalty0\vspace{1pt}
\begin{center}\small
\begin{tabularx}{\textwidth}{@{}L{3.0cm}XX@{}}
\toprule
\hrow \thead{Parameter} & \thead{BPSK} & \thead{QPSK} \\
\midrule
Bits per symbol & 1 & 2 (dibit) \\
Symbol period & $T_s = T_b$ & $T_s = 2T_b$ \\
Constellation & 2 points, $180^\circ$ apart & 4 points, $90^\circ$ apart \\
Null-to-null RF BW & $2R_b$ & $R_b$ \\
$\eta_B$ & 0.5 bps/Hz & \textbf{1 bps/Hz} \\
Bit error rate & $Q(\sqrt{2E_b/N_0})$ & $Q(\sqrt{2E_b/N_0})$, \textbf{the same} \\
Phase jumps & $180^\circ$ only & $0^\circ$, $\pm 90^\circ$ or $180^\circ$ \\
Circuit & one balance modulator & two, plus S/P and $90^\circ$ shifter \\
\bottomrule
\end{tabularx}
\end{center}
\penalty0\vspace{1pt}
{\color{sub}\footnotesize \mk{The one-line answer:} QPSK carries twice the data in the same
band at the same BER and the same $E_b/N_0$, paying only in circuit complexity.}

\Q{DPSK modulation: transmitter and receiver; PN sequence and why it is used}{\m{4+2+2}}
{\color{sub}\footnotesize \tP{1} \yr{72 Ma \m{4+2+2}} \gap$\cdot$\gap
\mk{the PN half of this question is really \S4.7 material, cross-refer it}}

\sbsr{0.47}{%
\lead{What DPSK is}
\begin{itemize}
  \item \textbf{Non-coherent form of PSK}: no synchronous carrier needed at the receiver.
  \item Information sits in the \mk{difference between the phases of two successive symbols},
        not in the absolute phase.
  \item Encoding rule: symbol \textbf{1} $=$ \textbf{no} phase transition, symbol \textbf{0}
        $=$ $180^\circ$ transition.
  \item So the present bit is decoded by comparing against the previous bit.
  \item $P_{e} = \tfrac{1}{2}e^{-E_b/N_0}$, about \textbf{3 dB worse} than coherent BPSK at
        high SNR.
\end{itemize}

\lead{Transmitter \m{2}}
\begin{enumerate}
  \item Binary sequence $m_k$ enters a \textbf{logic (differential encoder)} circuit.
  \item A \textbf{1-bit delay $T_b$} feeds the previous encoded bit $d_{k-1}$ back to it.
  \item Encoder output $d_k$ toggles when $m_k = 0$ and holds when $m_k = 1$.
  \item \textbf{Bipolar NRZ encoder} $\rightarrow$ $\pm 1$ levels.
  \item \textbf{Balance modulator} multiplies by the carrier $\rightarrow$ DPSK out.
\end{enumerate}}{%
\figT{c4_dbpsk_gen.png}
\penalty0\vspace{1pt}
{\color{sub}\footnotesize \textbf{DPSK transmitter}: encoder $+$ $T_b$ delay ahead of an
ordinary BPSK modulator. \yr{SST, Modulation extra note}}}

\penalty0\vspace{2pt}
\sbsr{0.50}{%
\figT{c4_dpsk_det.png}
\penalty0\vspace{1pt}
{\color{sub}\footnotesize \textbf{DPSK receiver}: the signal is its own phase reference.
\yr{SST, Modulation extra note}}}{%
\lead{Receiver \m{2}}
\begin{enumerate}
  \item \textbf{Bandpass filter} centred on $f_c$ removes out-of-band noise.
  \item Signal is multiplied by \mk{its own $T_b$-delayed copy} in the correlator.
  \item \textbf{Integrate and dump} over one bit gives an output proportional to
        $\cos\phi$, where $\phi$ is the phase difference between the two.
  \item $\phi = 0 \rightarrow$ output positive $\rightarrow$ decide \textbf{1}. \\
        $\phi = \pi \rightarrow$ output negative $\rightarrow$ decide \textbf{0}.
  \item \textbf{Decision device} at threshold $\lambda = 0$ V rebuilds the sequence.
\end{enumerate}
{\color{sub}\footnotesize \mk{No carrier recovery circuit anywhere in this diagram.} That is
the mark the examiner is looking for.}}

\penalty0\vspace{2pt}
\sbsr{0.56}{%
\figT{c4_dpsk_table.png}
\penalty0\vspace{1pt}
{\color{sub}\footnotesize Differential encoding worked through: the starred bit is the
arbitrary reference. Row 2 is $d_k$, row 3 its phase, last row the recovered data.
\yr{SST, Modulation extra note}}}{%
\lead{DPSK against BPSK}
\begin{itemize}
  \item \textbf{Bandwidth} $f_b$ against $2 f_b$ $\rightarrow$ DPSK halves it.
  \item \textbf{Carrier at demodulator}: not required against required.
  \item \textbf{Error probability}: higher against low.
  \item \textbf{Bit detection}: over two successive bit intervals against a single interval.
  \item \textbf{Circuit}: simpler receiver, but the encoder adds transmitter complexity.
\end{itemize}

\lead{(+) and (-)}
\begin{itemize}
  \item (+) No carrier recovery $\rightarrow$ cheap, small, fast to acquire.
  \item (+) Works where the channel rotates phase faster than a PLL can track.
  \item (-) $P_e$ roughly 3 dB worse than coherent BPSK.
  \item (-) \mk{Errors propagate in pairs}: one bad symbol corrupts the reference for the
        next.
\end{itemize}}

\penalty0\vspace{2pt}
\lead{PN sequence and why it is used \m{2+2} \; \small(full treatment in \S4.7)}
\begin{itemize}
  \item \textbf{PN (pseudo-noise) sequence} $=$ a coded sequence of 1s and 0s that looks
        random to anyone without the generator, but is \mk{completely deterministic and
        repeatable}.
  \item Generated by an $m$-stage shift register with XOR (parity) feedback; maximal-length
        period $= 2^m - 1$ chips.
  \item \textbf{Properties}: 1s and 0s roughly equiprobable; sharp autocorrelation peak and
        near-zero elsewhere; low cross-correlation between different codes; adding a shifted
        copy gives the same sequence in a new phase; easy to generate and synchronise.
  \item \textbf{Why used:} spreads the spectrum, so the signal hides under the noise floor
        (low probability of intercept) and shrugs off narrowband jamming.
  \item \textbf{Why used:} nearly orthogonal codes let many users share one band at once
        $\rightarrow$ CDMA.
  \item \textbf{Why used:} gives the receiver a unique pattern to correlate against, which is
        both the addressing scheme and the synchronisation reference.
\end{itemize}

\Q{OQPSK transmitter and receiver; why $\pi/4$-QPSK is preferred over OQPSK}{\m{5+2}}
{\color{sub}\footnotesize \tP{1} \yr{\textbf{71 Bh} \m{5+2}}}

\sbsr{0.53}{%
\lead{Offset QPSK (OQPSK), also called staggered QPSK}
\begin{itemize}
  \item \mk{Same circuit as QPSK, with one change}: the $Q$ arm is delayed by
        $T_b = T_s/2$ before its balance modulator.
  \item Effect: the even and odd streams can no longer change at the same instant.
  \item So \textbf{only one of I or Q switches at a time} $\rightarrow$ maximum phase change is
        $\pm 90^\circ$, never $180^\circ$.
  \item The envelope therefore \mk{never passes through zero}, and band-limiting produces
        only shallow dips instead of nulls.
  \item Hard-limiting or class-C amplification no longer regenerates the sidelobes, so OQPSK
        keeps its spectral occupancy in a real transmitter.
  \item Phase transitions happen every $T_b$ instead of every $T_s$, i.e. twice as often but
        half as violent.
  \item \textbf{PSD, BER and bandwidth are identical to QPSK.} The gain is purely in envelope
        behaviour.
\end{itemize}}{%
\figT{c4_oqpsk_wave.png}
\penalty0\vspace{1pt}
{\color{sub}\footnotesize \textbf{OQPSK time offset}: $m_I(t)$ changes at even multiples of
$T$, $m_Q(t)$ at odd multiples. Note the half-symbol offset. \yr{Rappaport Fig 5.30}}}

\sbsr{0.53}{%
\lead{Transmitter \m{2.5}}
\begin{enumerate}
  \item NRZ $\rightarrow$ serial-to-parallel into $m_I(t)$ and $m_Q(t)$.
  \item \mk{Insert a $T_b$ delay in the $Q$ branch.}
  \item $m_I$ $\times \cos 2\pi f_c t$, delayed $m_Q$ $\times \sin 2\pi f_c t$.
  \item Sum, then BPF $\rightarrow$ OQPSK out.
\end{enumerate}

\lead{Receiver \m{2.5}}
\begin{enumerate}
  \item Coherent, exactly the QPSK receiver: BPF, carrier recovery, two correlators,
        two decision circuits.
  \item \mk{The $T_b$ stagger is removed} by delaying the $I$ decision (or advancing $Q$)
        before the multiplexer, so the dibits re-align.
  \item Symbol timing recovery is taken from the transitions, which now occur every $T_b$.
\end{enumerate}}{%
\figT{c4_pi4_const.png}
\penalty0\vspace{1pt}
{\color{sub}\footnotesize \textbf{$\pi/4$-QPSK}: constellation (a) and constellation (b) are
used on alternate symbols, giving the 8-state pattern below with all allowed transitions
drawn. \yr{Rappaport Fig 5.31}}}

\penalty0\vspace{2pt}
\sbsr{0.50}{%
\lead{$\pi/4$-QPSK, the compromise}
\begin{itemize}
  \item Two QPSK constellations, offset from each other by $\pi/4$, used on
        \textbf{alternate symbols}.
  \item Signalling points therefore take \textbf{8} possible phase states, but only 4 are
        available in any one symbol.
  \item Maximum phase change is limited to $\pm 135^\circ$: worse than OQPSK's
        $90^\circ$, much better than QPSK's $180^\circ$.
  \item \mk{Every symbol carries a phase transition} (no $0^\circ$ option), which makes clock
        recovery easy.
  \item Usually sent \textbf{differentially encoded} as $\pi/4$-DQPSK.
  \item Used in IS-136 (US digital cellular), PDC and PHS (Japan), TETRA.
\end{itemize}}{%
\lead{\mk{Why $\pi/4$-QPSK is preferred over OQPSK} \m{2}}
\begin{enumerate}
  \item \textbf{It can be detected non-coherently.} Because it is differentially encoded, a
        differential or discriminator detector works. \mk{OQPSK must be detected coherently},
        and the half-symbol offset makes differential detection impossible.
  \item Non-coherent detection $\rightarrow$ \textbf{no carrier recovery circuit}
        $\rightarrow$ far simpler, cheaper, faster-acquiring receiver.
  \item \textbf{Better in multipath spread and fading}, where a coherent receiver cannot hold
        a phase reference. This is the mobile channel exactly.
  \item Guaranteed transition every symbol $\rightarrow$ reliable symbol timing recovery even
        on long runs of identical dibits.
  \item Cost accepted: envelope variation is larger than OQPSK's, so a mildly linear amplifier
        is still needed. \added
\end{enumerate}}

\hr

\T{4.4 Constant envelope modulation: BFSK, MSK, GMSK}

{\color{sub}\footnotesize \textbf{Constant envelope} (also called non-linear) means carrier
amplitude never changes, whatever the data. \mk{Why mobile radio loves it:} a class-C
amplifier can be run into saturation at 60 to 70\,\% efficiency instead of a linear amplifier
backed off to 20\,\%, which is the difference between a day and an afternoon of battery life.
Also: very little out-of-band spillage, simpler receiver, and high immunity to random FM noise
and Rayleigh fading.}

\sbsr{0.50}{%
\lead{BFSK, the starting point}
\begin{itemize}
  \item Two carrier frequencies $f_{c1}$ (mark) and $f_{c2}$ (space) carry 1 and 0.
  \item $s_i(t) = \sqrt{2E_b/T_b}\,\cos(2\pi f_{ci} t)$, $i = 1, 2$.
  \item Constant amplitude $\rightarrow$ constant envelope.
  \item Detected \textbf{coherently} (two correlators, compare outputs) or
        \textbf{non-coherently} (two BPFs $+$ envelope detectors $+$ comparator).
  \item (+) simple, insensitive to amplitude fluctuation, low $P_e$.
  \item (-) \mk{needs more bandwidth than ASK or PSK}.
  \item (-) If the two oscillators are free-running the phase jumps at every transition
        $\rightarrow$ spectral widening, large sidelobes.
  \item Fix: \textbf{CPFSK}, one oscillator whose frequency is shifted, so phase stays
        continuous.
\end{itemize}}{%
\lead{MSK $=$ the best CPFSK}
\begin{itemize}
  \item \mk{Definition 1 (as CPFSK):} a special case of CPFSK in which the peak frequency
        deviation is exactly \textbf{one quarter of the bit rate},
        $\Delta f = R_b/4$, giving modulation index $h = 2\Delta f\,T_b = 0.5$.
  \item $h = 0.5$ is the \mk{smallest index for which the two tones are orthogonal},
        $f_1 - f_2 = 1/(2T_b)$. Hence ``\textbf{minimum} shift''.
  \item Orthogonal $\rightarrow$ uncorrelated $\rightarrow$ the receiver separates them
        cleanly.
  \item \mk{Definition 2 (as OQPSK):} MSK is OQPSK in which the rectangular data pulses are
        replaced by \textbf{half-cycle sinusoidal} pulses.
  \item That pulse shaping is what makes the phase change \emph{linear} and bounded to
        $\pm \pi/2$ over a bit interval, so phase is continuous and smooth.
  \item Also called \textbf{Fast FSK}: the frequency shift is only half that of ordinary FSK.
  \item Good envelope, good spectral efficiency, good BER, no bandpass filter needed.
\end{itemize}}

\Q{MSK as a special type of CPFSK and of OQPSK; compare OQPSK, MSK and GMSK; which suits GSM}{\m{2}\m{4}\m{3+3+2}}
{\color{sub}\footnotesize \tF{4} \yr{\texttt{80 Ba}, \texttt{81 Ba} \m{3+3+2},
\textbf{80 Ch} \m{2}, \textbf{76 Bh} \m{4}} \gap$\cdot$\gap
\mk{MSK and GMSK together are asked in 6 of 22 papers}}

\penalty0\vspace{1pt}
\begin{center}\small
\begin{tabularx}{\textwidth}{@{}L{3.1cm}XXX@{}}
\toprule
\hrow \thead{} & \thead{OQPSK} & \thead{MSK} & \thead{GMSK} \\
\midrule
Family & linear & constant envelope & constant envelope \\
Pulse shape & rectangular & half-cycle sine & Gaussian \\
Null-to-null BW & $R_b$ & $1.5 R_b$ & below $R_b$ \\
99\,\% power BW & narrowest of the three & $1.20 R_b$ & $0.86 R_b$ at $BT = 0.25$ \\
Sidelobe roll-off & $f^{-2}$, slowest & $f^{-4}$ & fastest of the three \\
Envelope & near-constant & \textbf{constant} & \textbf{constant} \\
Power efficiency & $Q(\sqrt{2E_b/N_0})$ & the same & slight loss from ISI \\
Amplifier needed & mildly linear & class C, saturated & class C, saturated \\
Detection & coherent only & coherent or not & \textbf{FM discriminator works} \\
\bottomrule
\end{tabularx}
\end{center}
\penalty0\vspace{2pt}
{\color{sub}\footnotesize \mk{Read the table as a story}: OQPSK has the narrowest main lobe but
the slowest-decaying tails, MSK trades a wider main lobe for tails that die four times faster,
GMSK then squeezes the main lobe back down with a Gaussian filter and pays for it in ISI.}

\penalty0\vspace{2pt}
\sbsr{0.50}{%
\lead{\mk{Which is suitable for GSM, and why} \m{2}}
\begin{itemize}
  \item \textbf{GMSK with $BT = 0.3$.} Three reasons, in order of weight:
  \begin{enumerate}
    \item \textbf{Constant envelope} $\rightarrow$ the handset can use a saturated class-C
          power amplifier $\rightarrow$ battery life and cheap RF.
    \item \textbf{Compact spectrum} $\rightarrow$ fits the 200 kHz GSM carrier with acceptable
          adjacent channel interference. MSK alone leaks too much.
    \item \textbf{Detectable non-coherently} by a simple FM discriminator, so the receiver
          stays cheap.
  \end{enumerate}
  \item $BT = 0.3$ is the chosen compromise: tight enough spectrum, ISI penalty still small
        enough to be equalized out.
\end{itemize}}{%
\figT{c4_msk_psd.png}
\penalty0\vspace{1pt}
{\color{sub}\footnotesize PSD of MSK against GMSK at three $BT$ values, plus the occupied RF
bandwidth table in units of $R_b$. \mk{MSK's tails do not fall fast enough} to protect the
adjacent channel, which is the entire reason GMSK exists. \yr{SST Ch 4}}}

\Q{GMSK transmission and detection process, block diagram and constellation}{\m{6}\m{8}}
{\color{sub}\footnotesize \tF{4} \yr{\textbf{\texttt{79 Bh}} \m{2+6},
\texttt{81 Ba} \m{8}, \textbf{76 Bh} \m{4}, 70 Ma \m{8}} \gap$\cdot$\gap
\mk{79 Bh pairs it with ``advantages of digital modulation'', see \S4.1}}

\sbsr{0.50}{%
\lead{What GMSK is}
\begin{itemize}
  \item MSK preceded by a \textbf{premodulation Gaussian low-pass filter}.
  \item The Gaussian filter has \mk{no zero crossings and no overshoot}, so it smooths the
        rectangular data into a shape that suppresses the high-frequency content responsible
        for MSK's sidelobes.
  \item Transfer function
        $H(f) = e^{-\frac{\ln 2}{2}\left(f/B\right)^{2}}$,
        where $B$ is the \textbf{3 dB bandwidth} of the filter.
  \item The filter is fully described by $B$ and the symbol duration $T$, so GMSK is named by
        its \mk{$BT$ product}. $BT = \infty$ is plain MSK.
  \item \textbf{Smaller $BT$}: tighter spectrum, sidelobes fall away faster, but the pulse
        spreads over more symbols $\rightarrow$ more \mk{ISI} $\rightarrow$ higher BER.
  \item GSM: $BT = 0.3$. CDPD and Mobitex also use GMSK.
  \item Constellation: same 4 points as QPSK, but the trajectory between them
        \mk{follows the unit circle} instead of cutting across it, which is the constant
        envelope drawn.
\end{itemize}}{%
\figT{c4_gmsk_tx.png}
\penalty0\vspace{1pt}
{\color{sub}\footnotesize \textbf{GMSK transmitter, FM route.} \yr{SST Ch 4}}
\penalty0\vspace{4pt}
\figT{c4_gmsk_iq.png}
\penalty0\vspace{1pt}
{\color{sub}\footnotesize \textbf{GMSK transmitter, I-Q route.} \yr{SST, Fig 3.2.17}}}

\penalty0\vspace{2pt}
\sbsr{0.50}{%
\figT{c4_gmsk_rx.png}
\penalty0\vspace{1pt}
{\color{sub}\footnotesize \textbf{GMSK detector}, orthogonal coherent form. \yr{SST, Fig 3.2.18}}}{%
\lead{Transmission process \m{3}}
\begin{enumerate}
  \item Data $\rightarrow$ \textbf{bipolar NRZ encoder} $\rightarrow$ $\pm 1$ rectangular
        pulses.
  \item $\rightarrow$ \textbf{Gaussian LPF} of bandwidth $B$, which rounds every pulse edge.
  \item $\rightarrow$ \textbf{FM modulator} set to modulation index exactly 0.5, driven by the
        carrier oscillator $\cos\omega_c t$.
  \item Output is GMSK: constant amplitude, continuous smoothly varying phase.
  \item \textbf{I-Q alternative} (right-hand figure above): the filtered signal is integrated,
        split into in-phase and quadrature paths, each Gaussian-filtered and multiplied by
        $\cos\omega_c t$ / $\sin\omega_c t$, then subtracted. \mk{Advantage: the index stays at
        exactly 0.5 with no tuning or drift.}
\end{enumerate}

\lead{Detection process \m{3}}
\begin{enumerate}
  \item \textbf{Coherent (orthogonal) detector}, drawn left: GMSK splits into two arms.
  \item Each arm has a \textbf{balance modulator} fed by $\cos\omega_c t$ or $\sin\omega_c t$,
        then a \textbf{Gaussian filter} matched to the transmitter's.
  \item Each arm's \textbf{decision device} compares against its own threshold, giving
        $I(t)$ and $Q(t)$, which are recombined into the bit stream.
  \item \textbf{Non-coherent alternative}: a simple \textbf{FM discriminator} followed by an
        integrate-and-dump. Much cheaper, small BER penalty. This is what a GSM handset
        actually does.
\end{enumerate}}

\penalty0\vspace{2pt}
\lead{Worked example: sizing the Gaussian filter \added \; \small(tutorial problem, never a PYQ)}
\begin{asked}[PROBLEM]{SST tutorial \textperiodcentered{} not from any paper}
Find the 3-dB bandwidth for a Gaussian low pass filter used to produce 0.25 GMSK with a
channel data rate of $R_b = 270$ kbps. What is the 90\,\% power bandwidth in the RF channel?
\end{asked}

\sbsr{0.52}{%
\lead{Step 1: symbol duration}
\begin{itemize}
  \item Binary GMSK, so one bit per symbol and
        $T = 1/R_b = 1/270{,}000 = \mathbf{3.7\ \mu s}$.
\end{itemize}

\lead{Step 2: 3-dB bandwidth from the $BT$ product}
\begin{itemize}
  \item ``0.25 GMSK'' means $BT = 0.25$, so
        $B = 0.25/T = 0.25/3.7\ \mu\text{s} = \mathbf{67.57\ kHz}$.
\end{itemize}

\lead{Step 3: 90\,\% power RF bandwidth from the table}
\begin{itemize}
  \item Read the occupied RF bandwidth table (\S4.4 figure): row $BT = 0.25$, column 90\,\%
        gives $\mathbf{0.57 R_b}$.
  \item RF BW $= 0.57 \times 270{,}000 = \mathbf{153.9\ kHz}$.
\end{itemize}}{%
\lead{Answer}
\begin{itemize}
  \item 3-dB filter bandwidth $B = \mathbf{67.57\ kHz}$.
  \item 90\,\% power RF bandwidth $= \mathbf{153.9\ kHz}$.
\end{itemize}

\lead{Notes on this problem}
\begin{itemize}
  \item \mk{Real GSM uses $BT = 0.3$ and $R_b = 270.833$ kbps}, not the rounded values here.
        Answer with whatever the paper prints.
  \item The table is in units of $R_b$, so always multiply the tabulated figure by the bit
        rate. Forgetting that is the standard slip.
  \item Sanity check: 153.9 kHz sits inside GSM's 200 kHz carrier spacing, which is exactly
        the design intent.
\end{itemize}}

\hr

\T{4.5 M-ary modulation: MPSK, MFSK and QAM}

\sbsr{0.50}{%
\lead{The M-ary idea}
\begin{itemize}
  \item Send one of $M = 2^n$ distinct waveforms per signalling interval, carrying $n$ bits.
  \item Symbol duration $T_s = n T_b$, so the symbol rate drops by $n$ for the same bit rate.
  \item \mk{Bandwidth needed falls by a factor of $n$}, i.e. $\eta_B$ rises to $n$ bps/Hz
        in principle.
  \item Cost: the $M$ points crowd together, $d_{\min}$ shrinks, so \mk{more power is needed}
        for the same BER. $\eta_p$ falls.
  \item \textbf{MPSK}: all $M$ points on one circle of radius $\sqrt{E_s}$, differing only in
        phase. Constant envelope, but $d_{\min} = 2\sqrt{E_s}\sin(\pi/M)$ collapses as $M$
        grows.
  \item \textbf{MFSK}: $M$ orthogonal frequencies. Here bandwidth \emph{grows} with $M$ while
        power efficiency \emph{improves}: the opposite trade to MPSK.
  \item \textbf{QAM}: vary amplitude \emph{and} phase, so points fill a 2-D grid instead of a
        circle.
  \item 3G and later lean on these: EV-DO adapts between QPSK, 8-PSK and 16-QAM by channel
        quality.
\end{itemize}}{%
\figT{c4_qam16.png}
\penalty0\vspace{1pt}
{\color{sub}\footnotesize \textbf{16-QAM}, a $4\times4$ grid, 4 bits per symbol.
Nearest-neighbour distance is the same everywhere, which is why the grid beats the circle.
\yr{SST Ch 4}}}

\Q{Why do we need M-ary QAM? Compare 16-QAM and 16-PSK using constellation diagrams}{\m{2+6}}
{\color{sub}\footnotesize \tP{1} \yr{\textbf{77 Ch} \m{2+6}}}

\sbsr{0.46}{%
\lead{Why M-ary QAM is needed \m{2}}
\begin{itemize}
  \item Spectrum is licensed and finite, data demand is not. Binary schemes cap out near
        1 bps/Hz.
  \item M-ary raises $\eta_B$ to $\log_2 M$ bps/Hz without extra bandwidth.
  \item But \textbf{M-ary PSK alone fails at high $M$}: all points sit on one circle, so
        $d_{\min}$ falls like $\sin(\pi/M)$ and the power cost explodes.
  \item \mk{QAM uses the radial dimension too.} Spreading points over a 2-D grid gives a much
        larger $d_{\min}$ for the same average power.
  \item So M-ary QAM is the only practical way to reach 4, 6 or 8 bits per symbol on a
        band-limited channel.
\end{itemize}}{%
\figT{c4_16psk_qam.png}
\penalty0\vspace{1pt}
{\color{sub}\footnotesize \textbf{16-PSK} (left): 16 points on one circle, all at the same
power, spacing $22.5^\circ$. \textbf{16-QAM} (right): 16 points on a $4\times4$ grid, three
different amplitudes. Same 4 bits per symbol either way. \yr{SST Ch 4}}}

\penalty0\vspace{2pt}
\begin{center}\small
\begin{tabularx}{\textwidth}{@{}L{4.4cm}XX@{}}
\toprule
\hrow \thead{Parameter} & \thead{16-PSK} & \thead{16-QAM} \\
\midrule
Bits per symbol & 4 & 4 \\
Bandwidth & same & same (identical symbol rate) \\
Constellation & 16 points, one circle & $4\times4$ grid, 3 amplitude levels \\
Envelope & \textbf{constant} & varies \\
Amplifier & non-linear, efficient & \textbf{linear, backed off} \\
$d_{\min}$ at equal peak power & smaller & \textbf{larger} \\
Power for the same $d_{\min}$ & \textbf{1.6 dB more peak power} & reference \\
Power at equal $d_{\min}$ (average) & \textbf{4.19 dB more average power} \added & reference \\
Noise immunity & equal, if $d_{\min}$ is equalised & equal, at less power \\
\bottomrule
\end{tabularx}
\end{center}

\penalty0\vspace{2pt}
\sbs{%
\lead{The number the examiner wants}
\begin{itemize}
  \item \mk{At the same minimum Euclidean distance (the same noise immunity), 16-QAM needs
        1.6 dB less peak power than 16-PSK.}
  \item Measured on \emph{average} power instead, the gap is \textbf{4.19 dB} \added.
  \item Both figures describe the same fact: the grid packs 16 points more economically than
        the circle.
  \item \mk{Careful:} the lecture note's line ``noise immunity is same'' is only true when the
        two are compared at equal $d_{\min}$. At equal transmit power 16-QAM is clearly better.
\end{itemize}}{%
\lead{So why is 16-PSK used at all}
\begin{itemize}
  \item Constant envelope $\rightarrow$ the transmitter can saturate its amplifier.
  \item 16-QAM's amplitude variation forces a linear amplifier with 6 to 8 dB of back-off,
        which \mk{gives back much of the 4.19 dB}.
  \item Rule of thumb: QAM on the downlink (base station has mains power), constant envelope
        on the uplink (handset does not).
  \item Both need coherent detection and accurate gain control.
\end{itemize}}

\penalty0\vspace{2pt}
\sbsr{0.50}{%
\figT{c4_qam_mod.png}
\penalty0\vspace{1pt}
{\color{sub}\footnotesize \textbf{QAM modulator.} \yr{SST, Fig 3.2.21}}}{%
\lead{QAM modulator, walked through}
\begin{enumerate}
  \item Binary stream at $f_b$ enters a \textbf{2-bit serial-to-parallel converter}
        (4-bit for 16-QAM), clocked at the symbol rate $T_s$.
  \item Each half feeds a \textbf{D/A converter} that turns the bit group into a multi-level
        signal $d_1(t)$ and $d_2(t)$.
  \item $d_1(t)$ multiplies $\cos 2\pi f_c t$, $d_2(t)$ multiplies $\sin 2\pi f_c t$ (carrier
        through a $\pi/2$ shifter).
  \item Summer adds the two $\rightarrow$ a signal whose amplitude \emph{and} phase both carry
        data.
  \item \textbf{Demodulator} is the QPSK receiver with A/D converters instead of 1-bit
        decision devices, plus a carrier recovery circuit (raise to 4th power, BPF at $4f_c$,
        divide by 4).
\end{enumerate}}

\hr

\T{4.6 OFDM \; \small\mk{the single most examined topic in this chapter}}

{\color{sub}\footnotesize \textbf{Orthogonal frequency division multiplexing.} One high-rate
carrier replaced by many low-rate ones, spaced so they overlap without interfering. It is the
air interface of WiFi, WiMAX, LTE and 5G, and it is asked in \mk{9 of the 22 papers} here,
always with a block diagram attached.}

\Q{OFDM operation, block diagram of transmitter and receiver}{\m{4}\m{8}}
{\color{sub}\footnotesize \tS{9} \yr{\textbf{74 Bh} \m{8}, \textbf{\texttt{80 Bh}} \m{8},
\textbf{78 Ch} \m{8}, 71 Ma \m{8}, \textbf{70 Bh} \m{4}, \textbf{79 Ch} \m{4} (Tx only),
\textbf{80 Ch} \m{4} (Rx only), \textbf{\texttt{82 Bh}} \m{4+4}, 70 Ma \m{8}}
\gap$\cdot$\gap \mk{9 of 22 papers. Nothing else in the chapter comes close.}}

\sbsr{0.50}{%
\lead{The problem OFDM solves}
\begin{itemize}
  \item Push the bit rate up on a \textbf{single carrier} and the symbol period $T_s$ shrinks.
  \item Once $T_s$ approaches the channel's \textbf{delay spread} $\sigma_\tau$, echoes of
        symbol $n$ land on top of symbol $n{+}1$ $\rightarrow$ severe \textbf{ISI} and
        frequency-selective fading.
  \item Fixing that in the time domain needs an equalizer with a great many taps, which is
        expensive and must adapt continuously.
  \item \mk{OFDM's answer: do not make the symbol shorter, make it longer, and send many of
        them side by side.}
\end{itemize}}{%
\figT{c4_ofdm_fdm.png}
\penalty0\vspace{1pt}
{\color{sub}\footnotesize \mk{The figure that answers ``OFDM principle''.} FDM (top) wastes
spectrum on guard bands. OFDM (bottom) overlaps orthogonal subcarriers and needs none.
\yr{AS Sir, OFDM notes}}}

\sbsr{0.50}{%
\lead{The principle \m{4}}
\begin{itemize}
  \item Split the high-rate stream into \textbf{$N$ parallel low-rate substreams}.
  \item Each substream modulates its own \textbf{subcarrier}, using QPSK or QAM.
  \item Each subcarrier's symbol period becomes $N T_s$, now \mk{much longer than the delay
        spread} $\rightarrow$ each subchannel sees \textbf{flat} fading, not selective.
  \item Flat fading $\rightarrow$ one complex multiply per subcarrier undoes the channel.
        \mk{The many-tap equalizer disappears.}
  \item Subcarrier spacing is set to $\Delta f = 1/T$, which makes them
        \textbf{orthogonal}: each one's spectrum has a null exactly where every other one
        peaks.
  \item So the spectra may \mk{overlap without interfering}, and no guard bands are needed.
        Against FDM that is roughly a \textbf{50\,\% bandwidth saving}.
  \item The whole bank of $N$ modulators is computed as one \textbf{IFFT}, and the bank of
        demodulators as one \textbf{FFT}. No oscillator bank, no filter bank.
\end{itemize}}{%
\figT{c4_ofdm_concept.png}
\penalty0\vspace{1pt}
{\color{sub}\footnotesize Data is coded in the \textbf{frequency} domain, IFFT turns it into a
time waveform, the channel acts, FFT recovers each subcarrier separately.
\yr{AS Sir, OFDM notes}}}

\penalty0\vspace{2pt}
\sbsr{0.52}{%
\figT{c4_ofdm_block.png}
\penalty0\vspace{1pt}
{\color{sub}\footnotesize \textbf{Draw exactly this for the 8-mark question.} Transmitter on
top, receiver below, channel and noise between them. \yr{AS Sir, OFDM notes}}}{%
\lead{Transmitter chain \m{4}}
\begin{enumerate}
  \item \textbf{Modulation / symbol mapping}: bits $\rightarrow$ QPSK or QAM symbols (after
        channel coding and interleaving).
  \item \textbf{S/P}: serial symbol stream split into $N$ parallel subcarrier values, i.e. a
        frequency-domain block.
  \item \textbf{IFFT}: converts that block to $N$ time-domain samples. This \emph{is} the
        modulation.
  \item \textbf{P/S}: back to a serial sample stream.
  \item \textbf{Insert CP}: cyclic prefix prepended (see below).
  \item \textbf{D/A} $\rightarrow$ up-convert to RF $\rightarrow$ antenna.
\end{enumerate}

\lead{Receiver chain \m{4}}
\begin{enumerate}
  \item Down-convert $\rightarrow$ \textbf{A/D}.
  \item \textbf{Synchronise}: packet detection, then carrier frequency offset correction.
  \item \textbf{Remove CP}, discard the guard samples.
  \item \textbf{S/P} $\rightarrow$ block of $N$ time samples.
  \item \textbf{FFT} $\rightarrow$ back to the frequency domain, one value per subcarrier.
  \item \textbf{Equalize}: divide each subcarrier by its estimated $H[k]$, one tap each.
  \item \textbf{P/S} $\rightarrow$ \textbf{de-modulation / de-mapping} $\rightarrow$
        deinterleave and decode $\rightarrow$ data out.
\end{enumerate}
{\color{sub}\footnotesize \mk{Receiver is the transmitter read right to left.} Say that and you
have drawn half the answer.}}

\penalty0\vspace{2pt}
\sbsr{0.50}{%
\lead{Cyclic prefix, the detail that earns marks}
\begin{itemize}
  \item Even with long symbols, the tail of one OFDM symbol still smears into the next.
  \item A blank guard interval would work, but the hardware must transmit continuously and
        the exact delay spread is unknown.
  \item \mk{Solution: copy the last $T_g$ of the symbol and glue it on the front.}
  \item Two effects:
  \begin{itemize}
    \item Any echo arriving within $T_g$ falls inside the copied portion, so it corrupts only
          the prefix, which is thrown away $\rightarrow$ \textbf{ISI removed}.
    \item Linear convolution with the channel becomes \emph{circular} convolution, so after
          the FFT each subcarrier is simply scaled: $Y[k] = H[k]X[k]$ $\rightarrow$
          \textbf{one-tap equalization}.
  \end{itemize}
  \item Requirement: $T_g >$ maximum delay spread. In 802.11, CP\,:\,data $= 1:4$.
  \item Side benefit: the packet still decodes even if detected a little late.
  \item \mk{Cost: the CP is pure overhead}, up to 20 to 25\,\% of throughput.
\end{itemize}}{%
\figT{c4_cp.png}
\penalty0\vspace{1pt}
{\color{sub}\footnotesize The tail of the symbol is copied to the front as an extension.
\yr{AS Sir, OFDM notes}}

\lead{Two more practical details}
\begin{itemize}
  \item \textbf{Virtual (null) subcarriers} at the band edges are left unused: edge bins are
        the most vulnerable to frequency offset, and the gap suppresses out-of-band radiation.
  \item \textbf{Pilot subcarriers} carry known symbols so the receiver can track residual
        phase and estimate $H[k]$.
  \item 802.11a splits its 64 bins as \textbf{48 data $+$ 4 pilot $+$ 12 null} \added.
        The lecture note's ``48 of 64 occupied'' counts only the data bins.
\end{itemize}}

\creamq{OFDM: advantages and disadvantages}{\m{4}}
{\color{sub}\footnotesize \tP{1} \yr{\textbf{\texttt{82 Bh}} \m{4+4}} \gap$\cdot$\gap
\mk{asked together with ``draw the transmitting side'', so budget the time}}

\sbs{%
\lead{Advantages}
\begin{itemize}
  \item \textbf{High spectral efficiency}: overlapping orthogonal subcarriers, no guard bands,
        roughly twice FDM's.
  \item \textbf{Immune to ISI and frequency-selective fading}: long symbols plus the CP.
  \item \textbf{Simple equalization}: one complex multiply per subcarrier, no adaptive
        transversal filter.
  \item \textbf{Efficient to build}: FFT/IFFT in DSP, cost grows as $N\log N$, not $N^2$.
  \item \textbf{Robust to narrowband interference}: a jammer kills a few subcarriers, and
        coding $+$ interleaving recovers them.
  \item \textbf{Flexible}: power and modulation order can be set per subcarrier (adaptive bit
        loading); extends naturally to OFDMA and MIMO.
\end{itemize}}{%
\lead{Disadvantages}
\begin{itemize}
  \item \mk{High PAPR (peak-to-average power ratio)}: $N$ subcarriers can add in phase, so the
        amplifier needs heavy back-off $\rightarrow$ poor efficiency. \textbf{The main
        drawback.} It is why LTE uses SC-FDMA, not OFDMA, on the uplink.
  \item \textbf{Very sensitive to frequency offset and phase noise}: any error destroys
        orthogonality $\rightarrow$ inter-carrier interference (ICI).
  \item \textbf{Sensitive to Doppler}: a fast-moving mobile changes the channel within one
        long symbol.
  \item \textbf{CP overhead} wastes 20 to 25\,\% of both time and energy.
  \item \textbf{Tight synchronisation} needed in both time and frequency.
\end{itemize}
{\color{sub}\footnotesize \textbf{Used in}: DAB, DVB-T, 802.11a/g/n/ac, WiMAX, LTE downlink,
5G NR, ADSL.}}

\hr

\T{4.7 Spread spectrum: PN sequences, DSSS and FHSS}

\Q{Define spread spectrum modulation; DSSS with transmitter and receiver block diagrams}{\m{5}\m{6}\m{2+6}}
{\color{sub}\footnotesize \tF{4} \yr{\textbf{\texttt{82 Bh}} \m{2+6}, \texttt{81 Ba} \m{6},
\textbf{80 Ch} \m{2+5}, 73 Ma \m{4}} \gap$\cdot$\gap
\mk{spread spectrum overall: 7 of 22 papers}}

\sbsr{0.52}{%
\lead{Definition \m{2}}
\begin{itemize}
  \item A modulation in which the transmitted signal occupies a bandwidth \mk{far greater than
        the minimum needed to send the message}.
  \item Spreading is done by a \textbf{code independent of the data}, usually a PN sequence.
  \item The receiver \textbf{de-spreads} using a synchronised copy of the same code.
  \item \mk{Only a receiver holding the code can recover the message.} To anyone else the
        signal looks like a rise in the noise floor.
  \item \textbf{Processing gain}
        $PG = \dfrac{B_{ss}}{B} = \dfrac{T_b}{T_c} = N$ chips per bit. This number \emph{is}
        the jamming margin.
\end{itemize}

\lead{The two problems it answers}
\begin{itemize}
  \item How to secure information over a broadcast medium.
  \item How to survive deliberate jamming and external interference.
  \item (And, as a bonus, how to let many users share one band with no coordination.)
\end{itemize}}{%
\figT{c4_ss_model.png}
\penalty0\vspace{1pt}
{\color{sub}\footnotesize \textbf{General model of a spread spectrum system.} Note the
symmetry: the same pseudonoise generator appears at both ends, and it sits \emph{outside} the
channel coding. \yr{SST Ch 4, Fig 7.1}}}

\penalty0\vspace{2pt}
\sbsr{0.50}{%
\figT{c4_pn_gen.png}
\penalty0\vspace{1pt}
{\color{sub}\footnotesize \textbf{PN sequence generator}: $m$ D flip-flops in a shift register,
XOR parity feedback into $D_0$. Maximal length $2^m - 1$. \yr{SST Ch 4, Fig 12.3}}}{%
\lead{PN sequence properties \m{2}}
\begin{itemize}
  \item Coded sequence of 1s and 0s with \textbf{noise-like} statistics but deterministic
        generation.
  \item \textbf{Balance}: 1s and 0s occur with near-equal probability.
  \item \textbf{Shift property}: adding a shifted copy of a PN sequence gives back the same
        sequence in a different phase.
  \item \textbf{High autocorrelation} at zero lag, near zero elsewhere $\rightarrow$ the
        receiver can find the exact code phase.
  \item \textbf{Low cross-correlation} with other codes $\rightarrow$ other users look like
        noise.
  \item Easy to generate and to synchronise, in hardware, at chip rate.
\end{itemize}}

\penalty0\vspace{2pt}
\sbsr{0.50}{%
\lead{DSSS transmitter \m{3}}
\begin{enumerate}
  \item Original data bit stream at rate $R_b$, in $\pm 1$ form.
  \item \textbf{Chip (PN) generator} produces a code at rate $R_c = N R_b$, i.e. $N$ chips per
        data bit.
  \item \textbf{Multiplier} multiplies data by chips. \mk{Each data bit is replaced by $N$
        chips}, which is where the bandwidth expansion happens.
  \item The wideband product PSK-modulates the carrier $\rightarrow$ transmitted DSSS signal.
  \item Spectrum is now $N$ times wider and $N$ times lower in power density, so the signal
        can sit \emph{beneath} the noise floor.
\end{enumerate}

\lead{DSSS receiver \m{3}}
\begin{enumerate}
  \item BPF and PSK demodulation give back the wideband chip stream.
  \item A \textbf{local PN generator}, synchronised to the transmitter's, produces the same
        code.
  \item \textbf{Multiply again by the same code}: since $(\pm 1)^2 = +1$, the data collapses
        back to its original narrow bandwidth (de-spreading).
  \item \mk{Anything not multiplied by the matching code gets spread instead of de-spread},
        so interference and other users are pushed down by the processing gain.
  \item Narrowband filter $+$ integrate-and-dump $\rightarrow$ recovered data.
\end{enumerate}}{%
\figT{c4_dsss_block.png}
\penalty0\vspace{1pt}
{\color{sub}\footnotesize \textbf{DSSS spreading.} \yr{SST Ch 4}}
\penalty0\vspace{4pt}
\figT{c4_dsss_wave.png}
\penalty0\vspace{1pt}
{\color{sub}\footnotesize \textbf{DSSS waveform example}: one data bit, an 11-chip code, and
their product. Data 1 passes the code through, data 0 inverts it. \yr{SST Ch 4}}}

\penalty0\vspace{2pt}
\sbs{%
\lead{DSSS (+)}
\begin{itemize}
  \item Multiple access with no coordination $\rightarrow$ \textbf{CDMA}, IS-95 and 3G.
  \item Narrowband jamming is suppressed by the full processing gain.
  \item Multipath components arriving more than one chip apart are resolvable
        $\rightarrow$ \textbf{RAKE} receiver turns them into diversity (Ch 5).
  \item No frequency planning: every cell can use every frequency.
\end{itemize}}{%
\lead{DSSS (-)}
\begin{itemize}
  \item \mk{Near-far problem}: a close transmitter swamps a distant one, so tight and fast
        power control is mandatory.
  \item PN acquisition and tracking is slow and hardware-hungry.
  \item \textbf{Self-jamming}: code cross-correlation is low but not zero, so users interfere
        with each other as the load rises.
\end{itemize}}

\Q{FHSS operation with a neat block diagram}{\m{4}\m{5}}
{\color{sub}\footnotesize \tF{3} \yr{\textbf{79 Ch} \m{2+5}, \textbf{73 Bh} \m{4}, 73 Ma \m{4}}
\gap$\cdot$\gap \mk{as FHMA this is a TOP-tier Ch 7 question, see \S7.3}}

\sbsr{0.50}{%
\lead{Operation \m{4}}
\begin{itemize}
  \item FHSS uses \textbf{$M$ different carrier frequencies}, and the carrier hops among them
        in a pseudo-random pattern.
  \item Binary data is applied to an \textbf{M-ary FSK modulator}.
  \item A \textbf{pseudorandom code generator} drives a \textbf{frequency synthesizer}, which
        produces the whole range of hop frequencies from one reference.
  \item The synthesizer output at any instant is the \textbf{frequency hop}, chosen by the
        current $k$ code bits from a frequency table of $2^k$ entries.
  \item A \textbf{mixer} combines each hop with the MFSK signal $\rightarrow$ transmitted
        signal.
  \item \mk{The spectrum is spread sequentially, not instantaneously}: at any one moment the
        signal is narrowband, but over time it covers the whole band.
  \item The receiver runs the same code and synthesizer, so its local oscillator hops in
        lockstep and de-hops the signal back to a fixed IF.
\end{itemize}

\lead{Slow against fast hopping}
\begin{itemize}
  \item \textbf{Slow FHSS}: several symbols per hop ($R_h < R_s$). Cheaper synthesizer, used
        by GSM and Bluetooth.
  \item \textbf{Fast FHSS}: several hops per symbol ($R_h > R_s$). Each symbol gets frequency
        diversity, so a bad hop cannot kill it. Costs a fast synthesizer.
\end{itemize}}{%
\figT{c4_fhss_tx.png}
\penalty0\vspace{1pt}
{\color{sub}\footnotesize \textbf{FHSS transmitter.} Code generator $\rightarrow$ frequency
table $\rightarrow$ synthesizer $\rightarrow$ modulator. \yr{SST Ch 4}}
\penalty0\vspace{4pt}
\figT{c4_fhss_sel.png}
\penalty0\vspace{1pt}
{\color{sub}\footnotesize \textbf{Frequency selection}: each 3-bit group of the PN stream
indexes one row of the frequency table, which is the next hop. \yr{SST Ch 4}}}

\penalty0\vspace{2pt}
\sbs{%
\lead{FHSS (+)}
\begin{itemize}
  \item \mk{Immune to the near-far problem}: two users are almost never on the same frequency
        at the same instant, so no tight power control is needed.
  \item Resists partial-band jamming and frequency-selective fading, because the signal moves
        out of the bad band on the next hop.
  \item Synchronisation is easier and faster than DSSS's chip-level acquisition.
  \item Longer range for the same power, since instantaneous bandwidth stays narrow.
\end{itemize}}{%
\lead{FHSS (-)}
\begin{itemize}
  \item Needs a fast, agile and expensive \textbf{frequency synthesizer}.
  \item \textbf{Hits}: when two users pick the same hop, both symbols are lost, so error
        control coding is mandatory.
  \item Detection is usually non-coherent (the phase is lost at every hop), costing about
        3 dB.
  \item Total bandwidth is large, so it must be licensed or confined to an ISM band.
\end{itemize}
{\color{sub}\footnotesize \textbf{Bluetooth}: 79 hop channels, 1600 hops/s, 625 $\mu$s dwell.}}

\creamq{Advantages and major disadvantages of spread spectrum}{\m{2}\m{4}}
{\color{sub}\footnotesize \yr{\textbf{76 Bh} \m{4} (+), \textbf{77 Ch} \m{2} (-)}}

\sbs{%
\lead{Advantages \m{4}}
\begin{itemize}
  \item \textbf{Anti-jam}: a narrowband jammer is spread by $PG$ at the receiver, so its
        effective power drops by the processing gain.
  \item \textbf{Low probability of intercept}: the signal sits under the noise floor, so an
        eavesdropper cannot even tell it is there.
  \item \textbf{Privacy}: without the code the signal cannot be de-spread.
  \item \textbf{Multiple access}: near-orthogonal codes let many users occupy one band at once
        (CDMA).
  \item \textbf{Multipath rejection}: delayed copies decorrelate with the code and are either
        rejected or harvested by a RAKE receiver.
  \item \textbf{Graceful degradation}: adding users raises the noise floor a little rather than
        blocking calls outright.
  \item Mitigates ISI, and \mk{bandwidth efficiency actually rises as simultaneous users
        increase}.
\end{itemize}}{%
\lead{Major disadvantages \m{2}}
\begin{itemize}
  \item \mk{Large bandwidth} needed, many times the message bandwidth. Spectrum is expensive.
  \item Complex and costly transmitter and receiver: PN generation, acquisition, tracking.
  \item \textbf{Long acquisition time} before the first bit can be decoded.
  \item \textbf{Near-far problem} in DSSS $\rightarrow$ fast closed-loop power control needed.
  \item \textbf{Self-jamming} from non-zero code cross-correlation, which caps capacity.
  \item Synchronisation of the two PN generators must be exact to within a fraction of a chip.
\end{itemize}}

\hr

\T{4.8 Performance of modulation in fading channels \; \pillo{sub}{syllabus, no PYQ yet}}

\sbs{%
\lead{What the fading channel does to BER}
\begin{itemize}
  \item In \textbf{AWGN} the error probability falls \emph{exponentially} with $E_b/N_0$.
  \item In \textbf{flat Rayleigh fading} it falls only \emph{inversely}:
        $\bar{P_e} \approx 1/(4\,\overline{E_b/N_0})$ for coherent BPSK.
  \item Consequence: reaching $10^{-3}$ BER needs about \mk{7 dB in AWGN and about 24 dB in
        Rayleigh fading}. Tens of dB of extra power for the same quality. \added
  \item Cause: the deep fades dominate the average, not the typical signal level.
  \item Cure is \textbf{diversity}, not more power (Ch 5).
\end{itemize}}{%
\lead{The irreducible error floor}
\begin{itemize}
  \item In \textbf{frequency-selective} fading, ISI produces errors that
        \mk{no amount of transmit power removes}.
  \item The floor is set by the ratio $\sigma_\tau / T_s$, i.e. delay spread against symbol
        period.
  \item In \textbf{fast fading}, random FM from Doppler produces a second floor set by
        $f_d T_s$.
  \item Four ways out: equalization, spread spectrum $+$ RAKE, OFDM (long symbols), and
        diversity.
\end{itemize}}

\penalty0\vspace{2pt}
\sbs{%
\lead{Which scheme survives the mobile channel}
\begin{itemize}
  \item \textbf{Coherent} schemes suffer most: the receiver cannot track a phase reference
        that is being randomised by the channel.
  \item \textbf{Differential and non-coherent} schemes (DPSK, $\pi/4$-DQPSK, GMSK with a
        discriminator) degrade far more gracefully.
  \item \textbf{Constant envelope} schemes suit the handset, because the power amplifier can
        saturate.
  \item \textbf{Linear} schemes (QAM) need back-off and accurate gain control, so they belong
        on the downlink or on a fixed link.
\end{itemize}}{%
\lead{The trade-off table in one line each}
\begin{itemize}
  \item \textbf{BPSK}: most power efficient binary scheme, poorest $\eta_B$.
  \item \textbf{QPSK / OQPSK}: double $\eta_B$ at no BER cost, envelope problems.
  \item \textbf{MSK / GMSK}: constant envelope, non-coherent detection, GSM's choice.
  \item \textbf{M-QAM}: highest $\eta_B$, needs high SNR and a linear amplifier.
  \item \textbf{DSSS}: trades a huge bandwidth for interference immunity and multiple access.
  \item \textbf{OFDM}: turns a selective channel into many flat ones, pays in PAPR.
\end{itemize}}

\hr

\T{4.9 Last-minute recall for this chapter}

\sbs{%
\lead{Must memorise}
\begin{itemize}
  \item QPSK: 2 bits/symbol, $T_s = 2T_b$, null-null BW $= R_b$, \mk{same BER as BPSK}.
  \item MSK: CPFSK with $h = 0.5$, $\Delta f = R_b/4$, equals OQPSK with half-sine pulses.
  \item GMSK: MSK $+$ Gaussian LPF, $H(f) = e^{-\frac{\ln 2}{2}(f/B)^2}$,
        \mk{GSM uses $BT = 0.3$}.
  \item Smaller $BT$ $\rightarrow$ tighter spectrum, more ISI.
  \item OFDM: $\Delta f = 1/T$, IFFT at the transmitter, FFT at the receiver, CP $>$ delay
        spread, main drawback \mk{PAPR}.
  \item Spread spectrum: $PG = B_{ss}/B = T_b/T_c = N$.
  \item PN generator: $m$ stages, maximal length $2^m - 1$.
  \item 16-QAM against 16-PSK: \mk{1.6 dB less peak power} at equal $d_{\min}$, 4.19 dB less
        average power.
  \item $\pi/4$-QPSK: max phase change $135^\circ$, differentially detectable, IS-136 and PDC.
\end{itemize}}{%
\lead{Most repeated, in order}
\begin{enumerate}
  \item OFDM transmitter and receiver +Fig \tS{9}
  \item Spread spectrum, DSSS and FHSS +Fig \tS{7}
  \item MSK and GMSK, and why GSM chose GMSK \tS{6}
  \item QPSK transmission and detection +Fig \tF{4}
  \item OFDM (+) and (-) \tP{1}
  \item 16-QAM against 16-PSK \tP{1}
  \item DPSK transmitter and receiver, PN sequence \tP{1}
  \item OQPSK, and why $\pi/4$-QPSK beats it \tP{1}
\end{enumerate}
{\color{sub}\footnotesize \mk{Draw first, then annotate.} An unlabelled correct block diagram
still scores; three paragraphs with no diagram do not.}}

\penalty0\vspace{2pt}
{\color{sub}\footnotesize \mk{No paper in 2070--2082 has set a numerical from this chapter.}
The one calculation worth holding is the Gaussian-filter sizing in \S4.4, which is a tutorial
problem, not a PYQ. Everything examined here is descriptive, and the marks live in the figures.}
